\documentclass[12pt,oneside]{article}
\usepackage[margin=1in]{geometry}
\usepackage{amsmath}
\usepackage{amssymb}
\usepackage{mathtools}
\usepackage{xcolor}
\usepackage{fancybox}

\title{Full Algebraic Derivation of 2D Radiation Diffusion Implicit Stencil}
\author{}
\date{}

\newcommand{\labeleq}[1]{\tag{#1}}

\begin{document}

\maketitle

\tableofcontents
\newpage

% ============================================================
% Section 1: Continuous Equations
% ============================================================
\section{Continuous Equations in Cylindrical Coordinates}

We start with the coupled radiation and material energy equations in axisymmetric cylindrical coordinates $(r, z, t)$.

\subsection{Radiation Energy Equation}

The radiation energy density $E(r,z,t)$ evolves according to:
\begin{equation}
\frac{\partial E}{\partial t}
=
\frac{c}{3\sigma}\left(\frac{\partial^2 E}{\partial r^2} + \frac{1}{r}\frac{\partial E}{\partial r} + \frac{\partial^2 E}{\partial z^2}\right)
-
\chi c\sigma(E - U),
\label{eq:E_continuous}
\end{equation}
where
\begin{itemize}
\item $c$ is the speed of light,
\item $\sigma(T)$ is the temperature-dependent opacity (Rosseland mean), which varies in space via temperature,
\item $U(r,z,t)$ is the material energy density (to be coupled below),
\item the term $\frac{1}{r}\frac{\partial E}{\partial r}$ arises from cylindrical geometry.
\end{itemize}

\textbf{Important:} The form in Eq.~\eqref{eq:E_continuous} is written for notational simplicity. Since $\sigma(T(r,z))$ is \emph{space-dependent}, the operator $\frac{c}{3\sigma}\nabla^2$ is not linear in the usual sense. In discretization, we must use the \emph{conservative form} with diffusion coefficient $D(r,z) = \frac{c}{3\sigma(T(r,z))}$ varying at each location, giving:
\begin{equation}
\frac{\partial E}{\partial t}
=
\frac{1}{r}\frac{\partial}{\partial r}\left(rD\frac{\partial E}{\partial r}\right) + \frac{\partial}{\partial z}\left(D\frac{\partial E}{\partial z}\right)
-
\chi c\sigma(E - U).
\label{eq:E_continuous_conservative}
\end{equation}
This is the form we use throughout the derivation below.

\subsection{Material Energy Equation}

The material energy density $U(r,z,t)$ evolves according to:
\begin{equation}
\frac{1}{\beta}\,\frac{\partial U}{\partial t}
=
\chi c\sigma(E - U),
\label{eq:U_continuous}
\end{equation}
where $\beta = \frac{\partial U_R}{\partial U_m}$ is the ratio of radiation to material heat capacity (dimensionless coupling strength).

\subsection{Decomposition of the Diffusion Operator}

Define the space-dependent diffusion coefficient:
\begin{equation}
D(r,z,t)
\equiv
\frac{c}{3\sigma(T(r,z,t))}.
\end{equation}

The diffusion operator in conservative form (Eq.~\eqref{eq:E_continuous_conservative}) is:
\begin{equation}
\mathcal{L}[E]
\equiv
\frac{1}{r}\frac{\partial}{\partial r}\left(rD\frac{\partial E}{\partial r}\right) + \frac{\partial}{\partial z}\left(D\frac{\partial E}{\partial z}\right)
=
\mathcal{L}_r[E] + \mathcal{L}_z[E],
\end{equation}
where the radial part is:
\begin{equation}
\mathcal{L}_r[E]
\equiv
\frac{1}{r}\frac{\partial}{\partial r}\left(rD\frac{\partial E}{\partial r}\right)
\end{equation}
and the axial part is:
\begin{equation}
\mathcal{L}_z[E]
\equiv
\frac{\partial}{\partial z}\left(D\frac{\partial E}{\partial z}\right).
\end{equation}

In the numerical implementation, we discretize using a finite-volume approach where $D(r,z)$ is evaluated at cell faces and treated as piecewise constant over each control volume. The diffusion coefficients $D^n_{i,j\pm1/2}$ are then lagged in time (evaluated at $t^n$ from $U^n$).

The coupled equations in conservative form are:
\begin{align}
\frac{\partial E}{\partial t}
&=
\frac{1}{r}\frac{\partial}{\partial r}\left(rD\frac{\partial E}{\partial r}\right) + \frac{\partial}{\partial z}\left(D\frac{\partial E}{\partial z}\right)
-
\chi c\sigma(E - U),
\label{eq:E_decomposed}\\
\frac{1}{\beta}\,\frac{\partial U}{\partial t}
&=
\chi c\sigma(E - U).
\label{eq:U_decomposed}
\end{align}

% ============================================================
% Section 2: Time Discretization
% ============================================================
\section{Time Discretization via Backward Euler}

We discretize time using the backward (implicit) Euler scheme on a sequence of times
\[
t^n = t^{n-1} + \Delta t,
\]
where $\Delta t = \Delta t_n$ may vary with $n$.

\subsection{Implicit Euler for Radiation Energy}

Apply backward Euler to Eq.~\eqref{eq:E_decomposed}:
\begin{equation}
\frac{E^{n+1} - E^n}{\Delta t}
=
\frac{1}{r}\frac{\partial}{\partial r}\left(rD^n\frac{\partial E^{n+1}}{\partial r}\right) + \frac{\partial}{\partial z}\left(D^n\frac{\partial E^{n+1}}{\partial z}\right)
-
\chi c\sigma^n(E^{n+1} - U^{n+1}),
\end{equation}
where we use time-lagged space-dependent coefficients $D^n(r,z)$ and $\sigma^n(r,z)$ (computed from $U^n$) at the old time, while spatial derivatives of $E$ are evaluated at the new time (implicit in space, explicit in the coefficients).

Rearrange:
\begin{equation}
\frac{E^{n+1}}{\Delta t}
-
\frac{1}{r}\frac{\partial}{\partial r}\left(rD^n\frac{\partial E^{n+1}}{\partial r}\right) 
-
\frac{\partial}{\partial z}\left(D^n\frac{\partial E^{n+1}}{\partial z}\right)
+
\chi c\sigma^n E^{n+1}
-
\chi c\sigma^n U^{n+1}
=
\frac{E^n}{\Delta t}.
\end{equation}

\subsection{Implicit Euler for Material Energy}

\textbf{Note on LTE coupling:} The coupling factor $\chi$ (dimensionless) modulates the interaction strength between radiation and material in Local Thermodynamic Equilibrium (LTE). The factor appears in $A^n = \beta^n \Delta t \cdot \chi \cdot c\sigma^n$, with typical values $\chi = 1$ for full coupling or $\chi < 1$ for reduced coupling.

Apply backward Euler to Eq.~\eqref{eq:U_decomposed}, using time-lagged coefficients $\beta^n$ and $\sigma^n$ (computed from $U^n$):
\begin{equation}
\frac{1}{\beta^n}\,\frac{U^{n+1} - U^n}{\Delta t}
=
\chi c\sigma^n(E^{n+1} - U^{n+1}).
\end{equation}

Rearrange:
\begin{equation}
\frac{1}{\beta^n}\,U^{n+1}
+
\Delta t \cdot \chi \cdot c\sigma^n U^{n+1}
=
\frac{1}{\beta^n}\,U^n
+
\Delta t \cdot \chi \cdot c\sigma^n E^{n+1}.
\end{equation}

Factor:
\begin{equation}
\left(\frac{1}{\beta^n} + \Delta t \cdot \chi \cdot c\sigma^n\right) U^{n+1}
=
\frac{1}{\beta^n}\,U^n
+
\Delta t \cdot \chi \cdot c\sigma^n E^{n+1}.
\end{equation}

Define:
\begin{equation}
A^n
\equiv
\beta^n \Delta t \cdot \chi \cdot c\sigma^n.
\end{equation}
Then:
\begin{equation}
\left(1 + A^n\right) U^{n+1}
=
U^n
+
A^n E^{n+1}.
\end{equation}

Solve for $U^{n+1}$ in terms of $E^{n+1}$:
\begin{equation}
U^{n+1}
=
\frac{U^n + A^n E^{n+1}}{1 + A^n}
\label{eq:U_update}
\end{equation}

\subsection{Substituting the Material Update into the Radiation Equation}

Substitute Eq.~\eqref{eq:U_update} into the radiation equation:
\begin{equation}
\frac{E^{n+1}}{\Delta t}
-
\frac{1}{r}\frac{\partial}{\partial r}\left(rD^n\frac{\partial E^{n+1}}{\partial r}\right) 
-
\frac{\partial}{\partial z}\left(D^n\frac{\partial E^{n+1}}{\partial z}\right)
+
\chi c\sigma^n E^{n+1}
-
\chi c\sigma^n \frac{U^n + A^n E^{n+1}}{1 + A^n}
=
\frac{E^n}{\Delta t}.
\end{equation}

Expand the coupling term:
\begin{equation}
\chi c\sigma^n \frac{U^n + A^n E^{n+1}}{1 + A^n}
=
\frac{\chi c\sigma^n U^n}{1 + A^n}
+
\frac{\chi c\sigma^n A^n E^{n+1}}{1 + A^n}.
\end{equation}

Substitute back:
\begin{equation}
\frac{E^{n+1}}{\Delta t}
-
\frac{1}{r}\frac{\partial}{\partial r}\left(rD^n\frac{\partial E^{n+1}}{\partial r}\right) 
-
\frac{\partial}{\partial z}\left(D^n\frac{\partial E^{n+1}}{\partial z}\right)
+
\chi c\sigma^n E^{n+1}
-
\frac{\chi c\sigma^n A^n E^{n+1}}{1 + A^n}
=
\frac{E^n}{\Delta t}
+
\frac{\chi c\sigma^n U^n}{1 + A^n}.
\end{equation}

Combine the $E^{n+1}$ coefficients on the left:
\begin{equation}
\frac{1}{\Delta t} + \chi c\sigma^n - \frac{\chi c\sigma^n A^n}{1 + A^n}.
\end{equation}

Simplify the coupling diagonal:
\begin{align}
\chi c\sigma^n - \frac{\chi c\sigma^n A^n}{1 + A^n}
&=
\chi c\sigma^n\left(1 - \frac{A^n}{1 + A^n}\right)\\
&=
\chi c\sigma^n\left(\frac{1 + A^n - A^n}{1 + A^n}\right)\\
&=
\frac{\chi c\sigma^n}{1 + A^n}.
\end{align}

So the equation becomes:
\begin{equation}
\left(\frac{1}{\Delta t} + \frac{\chi c\sigma^n}{1 + A^n}\right) E^{n+1}
-
-
\frac{1}{r}\frac{\partial}{\partial r}\left(rD^n\frac{\partial E^{n+1}}{\partial r}\right) 
-
\frac{\partial}{\partial z}\left(D^n\frac{\partial E^{n+1}}{\partial z}\right)
=
\frac{E^n}{\Delta t}
+
\frac{\chi c\sigma^n}{1 + A^n} U^n.
\label{eq:E_implicit_final}
\end{equation}

Define the effective coupling coefficient:
\begin{equation}
C^n
\equiv
\frac{\chi c\sigma^n}{1 + A^n}.
\end{equation}

Then Eq.~\eqref{eq:E_implicit_final} reads:
\begin{equation}
\left(\frac{1}{\Delta t} + C^n\right) E^{n+1}
-
\frac{1}{r}\frac{\partial}{\partial r}\left(rD^n\frac{\partial E^{n+1}}{\partial r}\right) 
-
\frac{\partial}{\partial z}\left(D^n\frac{\partial E^{n+1}}{\partial z}\right)
=
\frac{E^n}{\Delta t}
+
C^n U^n.
\label{eq:E_final_coupled}
\end{equation}

% ============================================================
% Section 3: Spatial Discretization - Grid
% ============================================================
\section{Spatial Discretization}

\subsection{Grid Definition}

We discretize the domain into a 2D grid of cell-centers:
\begin{align}
z_i &= (i-1)\,\Delta z, \quad i = 0, 1, \ldots, N_z - 1, \quad \Delta z = \frac{L_z}{N_z - 1},\\
r_j &= \text{general radial nodes}, \quad j = 0, 1, \ldots, N_r - 1.
\end{align}

In general, $r_j$ may be non-uniform (e.g., uniform in the foam region and geometrically stretched in the gold layer). For simplicity in the derivation, we first assume uniform $\Delta r$, then note the modifications for non-uniform grids.

Stored fields are cell-centered:
\begin{equation}
E^{n}_{i,j} \approx E(r_j, z_i, t^n), \quad U^{n}_{i,j} \approx U(r_j, z_i, t^n).
\end{equation}

\subsection{Face Values and Averaging}

Spatial derivatives are computed at cell faces. The faces are located at:
\begin{align}
r_{j-\frac12} &= \frac{r_{j-1} + r_j}{2}, \quad \Delta r_j = r_j - r_{j-1},\\
z_{i-\frac12} &= \frac{z_{i-1} + z_i}{2}, \quad \Delta z = z_i - z_{i-1}.
\end{align}

For space-dependent coefficients like $D(T(r,z))$, we evaluate $T$ at cell centers to get $D_{i,j}$, then average $D$ to each face. Common averaging methods are:
\begin{align}
\text{Arithmetic: } \quad D_{i,j+\frac12} &\approx \frac{D_{i,j} + D_{i,j+1}}{2},\\
\text{Harmonic: } \quad D_{i,j+\frac12} &\approx \frac{2 D_{i,j} D_{i,j+1}}{D_{i,j} + D_{i,j+1}},\\
\text{Geometric: } \quad D_{i,j+\frac12} &\approx \sqrt{D_{i,j} D_{i,j+1}}.
\end{align}
The harmonic mean is often preferred for diffusion coefficients to ensure monotonicity, but the code supports all three. Similarly for z-faces. These face-averaged coefficients $D^n_{i,j\pm1/2}$ are then frozen at the old time $t^n$ during the implicit solve.

% ============================================================
% Section 4: Radial Diffusion Operator
% ============================================================
\section{Radial Diffusion Operator}

\subsection{Conservative Form}

The radial part of the diffusion operator with space-dependent $D(r,z)$ is:
\begin{equation}
\mathcal{L}_r[E]
=
\frac{1}{r}\frac{d}{dr}\left(rD\frac{\partial E}{\partial r}\right).
\end{equation}

This is the correct conservative form when $D$ varies with position. We discretize by defining radial fluxes at cell faces:
\begin{equation}
F^{r}_{i,j+\frac12}
\equiv
r_{j+\frac12}\,D^n_{i,j+\frac12}\,\frac{E^{n+1}_{i,j+1} - E^{n+1}_{i,j}}{\Delta r_{j+\frac12}},
\quad
\Delta r_{j+\frac12} = r_{j+1} - r_j.
\end{equation}

The divergence over a control volume at $(i,j)$ is then:
\begin{equation}
\left[\frac{1}{r}\frac{d}{dr}\left(rD\frac{\partial E}{\partial r}\right)\right]_{i,j}
\approx
\frac{1}{r_j \,\Delta r_j}\left(F^{r}_{i,j+\frac12} - F^{r}_{i,j-\frac12}\right),
\end{equation}
where $\Delta r_j$ is the width of the control volume around node $j$ (typically the average of $\Delta r_{j-\frac12}$ and $\Delta r_{j+\frac12}$).

\subsection{Discrete Form for Interior Nodes}

For interior nodes $1 \le j \le N_r - 2$, expanding the space-dependent operator:
\begin{equation}
\left[\frac{1}{r}\frac{d}{dr}\left(rD\frac{\partial E}{\partial r}\right)\right]_{i,j}
\approx
\frac{1}{r_j \,\Delta r_j}
\left[
r_{j+\frac12}\,D^n_{i,j+\frac12}\,\frac{E^{n+1}_{i,j+1} - E^{n+1}_{i,j}}{\Delta r_{j+\frac12}}
-
r_{j-\frac12}\,D^n_{i,j-\frac12}\,\frac{E^{n+1}_{i,j} - E^{n+1}_{i,j-1}}{\Delta r_{j-\frac12}}
\right].
\end{equation}

Here, $D^n_{i,j+\frac12}$ is the face-averaged diffusion coefficient at time $t^n$. Rearrange as a linear combination:
\begin{equation}
\approx
\frac{r_{j+\frac12}\,D^n_{i,j+\frac12}}{r_j\,\Delta r_j\,\Delta r_{j+\frac12}} E^{n+1}_{i,j+1}
-
\frac{r_{j+\frac12}\,D^n_{i,j+\frac12}}{r_j\,\Delta r_j\,\Delta r_{j+\frac12}} E^{n+1}_{i,j}
-
\frac{r_{j-\frac12}\,D^n_{i,j-\frac12}}{r_j\,\Delta r_j\,\Delta r_{j-\frac12}} E^{n+1}_{i,j}
+
\frac{r_{j-\frac12}\,D^n_{i,j-\frac12}}{r_j\,\Delta r_j\,\Delta r_{j-\frac12}} E^{n+1}_{i,j-1}.
\end{equation}

Define the effective cylindrical weights (with space-dependent coefficients):
\begin{align}
\lambda^{r,n}_{i,j+\frac12}
&\equiv
\frac{r_{j+\frac12}\,D^n_{i,j+\frac12}}{r_j\,\Delta r_j\,\Delta r_{j+\frac12}},\\
\lambda^{r,n}_{i,j-\frac12}
&\equiv
\frac{r_{j-\frac12}\,D^n_{i,j-\frac12}}{r_j\,\Delta r_j\,\Delta r_{j-\frac12}}.
\end{align}

Then the discrete radial contribution to the implicit step is:
\begin{equation}
\left[\frac{1}{r}\frac{d}{dr}\left(rD^n\frac{\partial E}{\partial r}\right)\right]_{i,j}^{n+1}
\approx
\lambda^{r,n}_{i,j+\frac12} E^{n+1}_{i,j+1}
-
\left(\lambda^{r,n}_{i,j+\frac12} + \lambda^{r,n}_{i,j-\frac12}\right) E^{n+1}_{i,j}
+
\lambda^{r,n}_{i,j-\frac12} E^{n+1}_{i,j-1}.
\label{eq:radial_stencil}
\end{equation}

% ============================================================
% Section 5: Axial Diffusion Operator
% ============================================================
\section{Axial Diffusion Operator}

\subsection{Continuous Form}

The axial part with space-dependent $D(r,z)$:
\begin{equation}
\mathcal{L}_z[E]
=
\frac{\partial}{\partial z}\left(D\frac{\partial E}{\partial z}\right).
\end{equation}

\subsection{Discrete Form for Interior Nodes}

For interior axial nodes $i = 1, \ldots, N_z - 2$, the space-dependent operator is approximated as:
\begin{equation}
\left[\frac{\partial}{\partial z}\left(D\frac{\partial E}{\partial z}\right)\right]_{i,j}
\approx
\frac{1}{(\Delta z)^2}\left[D^n_{i+1,j}\,E^{n+1}_{i+1,j} - (D^n_{i+1,j} + D^n_{i,j})\,E^{n+1}_{i,j} + D^n_{i,j}\,E^{n+1}_{i-1,j}\right].
\end{equation}

However, for simplicity and to maintain stability, a common approximation averages the diffusion coefficient to the cell center:
\begin{equation}
D^n_{i,j} \approx \frac{D^n_{i+1,j} + D^n_{i,j} + D^n_{i-1,j}}{3},
\end{equation}
which then gives the simpler form:
\begin{equation}
\left[\frac{\partial}{\partial z}\left(D\frac{\partial E}{\partial z}\right)\right]_{i,j}
\approx
D^n_{i,j}\,\frac{E^{n+1}_{i+1,j} - 2E^{n+1}_{i,j} + E^{n+1}_{i-1,j}}{(\Delta z)^2}.
\end{equation}

With implicit time stepping:
\begin{equation}
\left[\frac{\partial}{\partial z}\left(D^n\frac{\partial E}{\partial z}\right)\right]_{i,j}^{n+1}
\approx
\frac{D^n_{i,j}}{(\Delta z)^2}
\left[E^{n+1}_{i+1,j} - 2E^{n+1}_{i,j} + E^{n+1}_{i-1,j}\right].
\end{equation}

Define the axial weight:
\begin{equation}
\lambda^{z,n}_{i,j}
\equiv
\frac{D^n_{i,j}}{(\Delta z)^2}.
\end{equation}

Then:
\begin{equation}
\left[D\frac{\partial^2 E}{\partial z^2}\right]_{i,j}^{n+1}
\approx
\lambda^{z,n}_{i,j} E^{n+1}_{i+1,j}
-
2\lambda^{z,n}_{i,j} E^{n+1}_{i,j}
+
\lambda^{z,n}_{i,j} E^{n+1}_{i-1,j}.
\label{eq:axial_stencil}
\end{equation}

% ============================================================
% Section 6: Complete Interior Stencil
% ============================================================
\section{Complete Interior 5-Point Stencil}

\subsection{Combining All Terms}

Start from Eq.~\eqref{eq:E_final_coupled} and apply time discretization:
\begin{equation}
\left(\frac{1}{\Delta t} + C^n\right) E^{n+1}_{i,j}
-
\left[\mathcal{L}_r + \mathcal{L}_z\right] E^{n+1}_{i,j}
=
\frac{E^n_{i,j}}{\Delta t}
+
C^n U^n_{i,j}.
\end{equation}

Substitute the discrete forms from Eqs.~\eqref{eq:radial_stencil} and \eqref{eq:axial_stencil}:
\begin{equation}
\left(\frac{1}{\Delta t} + C^n\right) E^{n+1}_{i,j}
-\left[
\lambda^{r,n}_{i,j+\frac12} E^{n+1}_{i,j+1}
-
(\lambda^{r,n}_{i,j+\frac12} + \lambda^{r,n}_{i,j-\frac12}) E^{n+1}_{i,j}
+
\lambda^{r,n}_{i,j-\frac12} E^{n+1}_{i,j-1}
\right]
\end{equation}
\begin{equation}
-\left[
\lambda^{z,n}_{i,j} E^{n+1}_{i+1,j}
-
2\lambda^{z,n}_{i,j} E^{n+1}_{i,j}
+
\lambda^{z,n}_{i,j} E^{n+1}_{i-1,j}
\right]
=
\frac{E^n_{i,j}}{\Delta t}
+
C^n U^n_{i,j}.
\end{equation}

Collect all terms on the left:
\begin{align}
&\left(\frac{1}{\Delta t} + C^n + \lambda^{r,n}_{i,j+\frac12} + \lambda^{r,n}_{i,j-\frac12} + 2\lambda^{z,n}_{i,j}\right) E^{n+1}_{i,j}\\
&\quad - \lambda^{r,n}_{i,j+\frac12} E^{n+1}_{i,j+1}
- \lambda^{r,n}_{i,j-\frac12} E^{n+1}_{i,j-1}
- \lambda^{z,n}_{i,j} E^{n+1}_{i+1,j}
- \lambda^{z,n}_{i,j} E^{n+1}_{i-1,j}\\
&=
\frac{E^n_{i,j}}{\Delta t}
+
C^n U^n_{i,j}.
\end{align}

Rearrange into standard form:
\begin{equation}
\boxed{
\begin{alignedat}{2}
&\text{diagonal coeff: }&&
\frac{1}{\Delta t} + C^n + \lambda^{r,n}_{i,j+\frac12} + \lambda^{r,n}_{i,j-\frac12} + 2\lambda^{z,n}_{i,j},\\
&\text{r-neighbors: }&&
-\lambda^{r,n}_{i,j\pm\frac12},\\
&\text{z-neighbors: }&&
-\lambda^{z,n}_{i,j},\\
&\text{RHS: }&&
\frac{E^n_{i,j}}{\Delta t}
+
C^n U^n_{i,j}.
\end{alignedat}
}
\label{eq:5point_interior}
\end{equation}

This is the full interior 5-point stencil (or 4-point in 1D) used in the implicit solver.

% ============================================================
% Section 7: Boundary Conditions
% ============================================================
\section{Boundary Conditions}

\subsection{Axial Boundary Conditions}

\subsubsection{Dirichlet Boundary at $z = 0$ (Drive)}

At the irradiated surface, we impose:
\begin{equation}
E^{n+1}_{0,j} = E_{\mathrm{drive}}(t^{n+1}, j),
\end{equation}
where $E_{\mathrm{drive}}(t)$ is computed from the imposed drive temperature $T_D(t)$ via $E_{\mathrm{drive}} = a T_D^4$ with $a = 4\sigma_{\mathrm{SB}}/c$.

This node is eliminated from the linear system; instead, its value is imposed after the solve.

\subsubsection{Marshak Boundary at $z = 0$}

Alternatively, the Marshak boundary condition relates the boundary value to the outward flux:
\begin{equation}
a T_D^4(t) = E(0,j,t) - \frac{2D(0,j,t)}{c}\,\left.\frac{\partial E}{\partial z}\right|_{z=0}.
\end{equation}

Using backward differencing:
\begin{equation}
\left.\frac{\partial E}{\partial z}\right|_{z=0}
\approx
\frac{E_{1,j}^{n+1} - E_{0,j}^{n+1}}{\Delta z}.
\end{equation}

Substituting:
\begin{equation}
a T_D^4(t^{n+1})
=
E_{0,j}^{n+1} - \frac{2D^n_{0,j}}{c}\,\frac{E_{1,j}^{n+1} - E_{0,j}^{n+1}}{\Delta z}.
\end{equation}

Multiply by $\Delta z$ and rearrange:
\begin{equation}
a T_D^4(t^{n+1}) \Delta z
=
E_{0,j}^{n+1} \Delta z - \frac{2D^n_{0,j}}{c}\,(E_{1,j}^{n+1} - E_{0,j}^{n+1}).
\end{equation}

\begin{equation}
a T_D^4(t^{n+1}) \Delta z
=
E_{0,j}^{n+1}\left(\Delta z + \frac{2D^n_{0,j}}{c}\right) - \frac{2D^n_{0,j}}{c}\,E_{1,j}^{n+1}.
\end{equation}

Define $\alpha_j = \frac{2D^n_{0,j}}{c\,\Delta z}$. Then:
\begin{equation}
a T_D^4(t^{n+1}) = (1 + \alpha_j)\,E_{0,j}^{n+1} - \alpha_j\,E_{1,j}^{n+1}.
\label{eq:marshak_row_z0}
\end{equation}

This is the discrete Marshak row at $z=0$, replacing the interior 5-point stencil for that location.

\subsubsection{Dirichlet Boundary at $z = L_z$ (Bath)}

At the far end:
\begin{equation}
E^{n+1}_{N_z-1,j} = E_{\mathrm{bath}} = a(300\text{ K})^4.
\end{equation}

\subsection{Radial Boundary Conditions}

\subsubsection{Axis Symmetry at $r = 0$}

By cylindrical symmetry:
\begin{equation}
\left.\frac{\partial E}{\partial r}\right|_{r=0} = 0.
\end{equation}

This is enforced by setting $E_{i,0} = E_{i,1}$, which removes the $1/r$ singularity.

\subsubsection{Dirichlet Boundary at $r = R_{\mathrm{tot}}$ (Fixed Bath)}

Impose:
\begin{equation}
E^{n+1}_{i,N_r-1} = E_{\mathrm{bath}}.
\end{equation}

\subsubsection{Neumann Boundary at $r = R_{\mathrm{tot}}$ (No Flux)}

Impose:
\begin{equation}
\left.\frac{\partial E}{\partial r}\right|_{r=R_{\mathrm{tot}}} = 0,
\end{equation}
which eliminates the outer-face flux contribution in the stencil.

\subsubsection{Marshak Boundary at $r = R_{\mathrm{tot}}$ (Radiating Wall)}

The Marshak condition at the outer surface is:
\begin{equation}
a T_{\mathrm{wall}}^4(t) = E(N_r-1,i,t) - \frac{2D(N_r-1,i,t)}{c}\,\left.\frac{\partial E}{\partial r}\right|_{r=R_{\mathrm{tot}}}.
\end{equation}

Using backward differencing:
\begin{equation}
\left.\frac{\partial E}{\partial r}\right|_{r=R_{\mathrm{tot}}}
\approx
\frac{E_{i,N_r-1}^{n+1} - E_{i,N_r-2}^{n+1}}{\Delta r_{\mathrm{outer}}}.
\end{equation}

Substituting and rearranging similarly to Eq.~\eqref{eq:marshak_row_z0}:
\begin{equation}
\boxed{
a T_{\mathrm{wall}}^4(t^{n+1})
=
(1 + \alpha_i)\,E_{i,N_r-1}^{n+1} - \alpha_i\,E_{i,N_r-2}^{n+1},
}
\label{eq:marshak_row_r_outer}
\end{equation}
where $\alpha_i = \frac{2D^n_{i,N_r-1}}{c\,\Delta r_{\mathrm{outer}}}$.

This row is included in the linear system to replace the interior 5-point stencil at the outer radial boundary.

% ============================================================
% Section 8: Linear System Assembly
% ============================================================
\section{Linear System Assembly}

\subsection{Unknown Vector Organization}

The implicit solve solves for $E^{n+1}$ on an interior subset of the grid. Typically:
\begin{itemize}
\item If Marshak boundary at $z=0$: unknowns are $E^{n+1}_{i,j}$ for $i=0,\ldots,N_z-2$ and $j=0,\ldots,N_r-1$.
\item If Dirichlet at $z=0$: unknowns are $E^{n+1}_{i,j}$ for $i=1,\ldots,N_z-2$ and $j=0,\ldots,N_r-1$.
\item Boundary nodes at $z=L_z$ and possibly $r=R$ are either imposed after the solve or handled specially.
\end{itemize}

The total number of unknowns is:
\begin{equation}
n_{\text{unknown}} = n_z \times n_r,
\end{equation}
where $n_z$ and $n_r$ are the counts of unknown rows in each direction.

\subsection{Row-Wise Assembly}

For each unknown $(i,j)$, the corresponding matrix row contains:
\begin{itemize}
\item Diagonal entry: coefficient of $E^{n+1}_{i,j}$,
\item Off-diagonal entries: coefficients of neighboring $E^{n+1}_{i\pm1,j}$ and $E^{n+1}_{i,j\pm1}$,
\item RHS: constant vector entries.
\end{itemize}

The sparsity pattern is structured: each row has at most 5 nonzeros (2D grid), or fewer at boundaries.

\subsection{CSR Storage}

To store the matrix efficiently in Compressed Sparse Row (CSR) format:
\begin{align}
&\text{indptr}: \text{ row start pointers (length } n_{\text{unknown}} + 1),\\
&\text{indices}: \text{ column indices of stored values,}\\
&\text{data}: \text{ numeric coefficients.}
\end{align}

For each row, we store only the nonzero entries in sorted column order, reducing memory from $O(n_{\text{unknown}}^2)$ dense to $O(n_{\text{unknown}})$ sparse.

% ============================================================
% Section 9: Material Update
% ============================================================
\section{Material Energy Update}

After solving for $E^{n+1}$, update $U^{n+1}$ using Eq.~\eqref{eq:U_update}:
\begin{equation}
U^{n+1}_{i,j}
=
\frac{U^n_{i,j} + A^n_{i,j} E^{n+1}_{i,j}}{1 + A^n_{i,j}},
\end{equation}
where $A^n_{i,j} = \beta^n_{i,j} \Delta t \cdot \chi \cdot c\sigma^n_{i,j}$.

This is a local (cell-by-cell) update, not requiring a linear solve.

% ============================================================
% Section 10: Self-Similar Model Integration
% ============================================================
\section{Self-Similar Model Integration}

In the self-similar model, the coupling parameter $\beta(T)$ (hence $\beta^n$ at the old time) and the heat capacity are related to material parameters via:
\begin{equation}
\beta(T) = \text{Cv}_R(T) / \text{Cv}_m(T),
\end{equation}
which is the ratio of radiation to material heat capacity (dimensionless). This is temperature-dependent and must be evaluated at each location from the old material energy $U^n$ to get $\beta^n_{i,j}$.

The opacity and heat capacity are functions of temperature:
\begin{align}
\sigma(T) &= \frac{1}{g T^\alpha \rho^{-\lambda - 1}},\\
\text{Cv}_m(T) &= f \beta_{\exp} T^{\beta_{\exp} - 1} \rho^{-\mu + 1},\\
\text{Cv}_R(T) &= 4 a T^3.
\end{align}

At each timestep, $T^n = (U^n / a)^{1/4}$ is computed from the old material energy, and all temperature-dependent quantities ($D^n, \sigma^n, C^n$) are evaluated there.

% ============================================================
% Appendix: Summary
% ============================================================
\appendix

\section{Quick Reference: Final Stencil Summary}

\subsection{Interior 5-Point Stencil}

For interior nodes $(i,j)$ with $1 \le i \le N_z-2$ and $1 \le j \le N_r-2$ (or adjusted ranges for Marshak):

\subsubsection{Matrix Coefficients}

\begin{align}
a_{\text{self}} &= \frac{1}{\Delta t} + C^n + \lambda^{r,n}_{i,j+\frac12} + \lambda^{r,n}_{i,j-\frac12} + 2\lambda^{z,n}_{i,j},\\
a_{i+1,j} &= -\lambda^{z,n}_{i,j},\\
a_{i-1,j} &= -\lambda^{z,n}_{i,j},\\
a_{i,j+1} &= -\lambda^{r,n}_{i,j+\frac12},\\
a_{i,j-1} &= -\lambda^{r,n}_{i,j-\frac12}.
\end{align}

\subsubsection{RHS}

\begin{equation}
b_{i,j} = \frac{E^n_{i,j}}{\Delta t} + C^n_{i,j} U^n_{i,j},
\end{equation}
where $C^n_{i,j} = \frac{c\sigma^n_{i,j}}{1 + A^n_{i,j}}$ and $A^n_{i,j} = \beta^n_{i,j} \Delta t \cdot \chi \cdot c\sigma^n_{i,j}$.

\subsection{Marshak Rows at $z=0$}

\begin{align}
a_{\text{self}} &= 1 + \alpha_j, \quad \alpha_j = \frac{2D^n_{0,j}}{c\,\Delta z},\\
a_{1,j} &= -\alpha_j,\\
b_{0,j} &= a T_D^4(t^{n+1}).
\end{align}

\subsection{Marshak Rows at $r=R_{\mathrm{tot}}$}

\begin{align}
a_{\text{self}} &= 1 + \alpha_i, \quad \alpha_i = \frac{2D^n_{i,N_r-1}}{c\,\Delta r_{\mathrm{outer}}},\\
a_{i,N_r-2} &= -\alpha_i,\\
b_{i,N_r-1} &= a T_{\mathrm{wall}}^4(t^{n+1}).
\end{align}

\end{document}
